\section{Finite returns and the Bloch coefficient}\label{sec:ideal}

We identify the finite return terms using the walk with every slot
available. Its site motion uses all interaction bonds, so its diffusion
matrix is \(D_{S,J}\), rather than a matrix determined by the three
bonds chosen for the geometric estimates. The comparison has two steps:
a coupling makes the expected mass error \(o_K(1)\), and splitting a
return at its midpoint multiplies this error by the diffusion bound
\(O(\beta^{-3/2})\).

Fix \(K\), \(E\), \(i\notin E\), \(P=E\cup\{i\}\), and
\(R=R_K\) as in Section~\ref{sec:returns}. Write \(x\) for the site
of \(i\) and \(z=\bar x\). For a direction word
\(\mathbf d=(d_1,\ldots,d_n)\in\{+,-\}^n\), traversed with
\(d_1\) first, define
\[
 W(\mathbf d)=U_{d_n}\cdots U_{d_1},\qquad
 W_R(\mathbf d)=U_{d_n,R}\cdots U_{d_1,R}.
\]
The factors are formed in one fixed exposure; each traversal then uses
fresh transition randomness. The ideal kernels have all slots available,
exchange rate \(J(y-x)/2\) on each slot edge, and fresh uniform site
permutations at every seam. The reverse factors use the inverse
permutations and the reversed order of operations. Denote the full
ideal word over \(\mathbb Z^3\), with no killing, by
\(\widehat W_\infty(\mathbf d)\). Denote its version killed outside
\(\Lambda^\circ\) and outside \(\mathfrak B_R(z)\) by
\(\widehat W_R(\mathbf d)\). Ideal and exposed kernels assign each
instantaneous operation to the same endpoint and use the same matching
layers. The ideal kernels are deterministic transition matrices, with
their fresh randomness already averaged.

\subsection{The full interaction kernel}

The site of a full ideal path jumps by \(v\) at rate
\(\ell J(v)/2=SJ(v)\). Its total jump rate is the finite constant
\(\lambda_J=S\sum_vJ(v)\). Neither a seam nor a reversal changes
this site generator. With
\[
 \phi(k)=S\sum_vJ(v)(1-\cos(k\cdot v)),\qquad
 \kappa=\frac1{(4\pi)^{3/2}\sqrt{\det D_{S,J}}},
\]
the site transition kernel is
\begin{equation}\label{eq:ideal-fourier}
 p_t(y,w)=\int_{[-\pi,\pi]^3}
     e^{\mathrm i k\cdot(w-y)}e^{-t\phi(k)}\,
                         \frac{dk}{(2\pi)^3}.
\end{equation}

\begin{lemma}[Ideal return coefficient]\label{lem:ideal-kernel}
The full site kernel satisfies
\begin{equation}\label{eq:ideal-site-bounds}
 \sup_{y,w}p_t(y,w)\le C(1+t)^{-3/2},\qquad
 p_t(x,x)=t^{-3/2}(\kappa+o(1)).
\end{equation}
The first estimate holds for all \(t\ge0\); the second is as
\(t\to\infty\). For every fixed \(K\), every \(1\le n\le2K\), and every direction
word of length \(n\),
\begin{equation}\label{eq:ideal-diagonal}
 (\widehat W_\infty(\mathbf d))_{ii}
 =\ell^{-1}p_{n\beta}(x,x)
 =\beta^{-3/2}\bigl((\kappa/\ell)n^{-3/2}+o_K(1)\bigr).
\end{equation}
The error is uniform in the word, cut, and initial slot. Always,
\begin{equation}\label{eq:ideal-killed-upper}
 0\le(\widehat W_R(\mathbf d))_{ii}
 \le(\widehat W_\infty(\mathbf d))_{ii}.
\end{equation}
If \(\mathfrak B_R(z)\subseteq\Lambda^\circ\), then
\begin{equation}\label{eq:ideal-killed-interior}
 0\le(\widehat W_\infty(\mathbf d))_{ii}
          -(\widehat W_R(\mathbf d))_{ii}
 \le C_K e^{-c\beta}.
\end{equation}
\end{lemma}

\begin{proof}
At a zero of \(\phi\), each \(k\cdot v\) with \(J(v)>0\)
belongs to \(2\pi\mathbb Z\). Since these \(v\) generate
\(\mathbb Z^3\), this forces \(k\in2\pi\mathbb Z^3\). Thus zero
is the only zero in the integration cube. Moreover,
\[
 \phi(k)=k^{\mathsf T}D_{S,J}k+O(|k|^4).
\]
The matrix is positive definite because the support spans
\(\mathbb R^3\). Compactness away from zero now gives
\(\phi(k)\ge c|k|^2\) throughout \([-\pi,\pi]^3\).
Taking absolute values in \eqref{eq:ideal-fourier} proves the heat
kernel upper bound, including \(t\le1\) by enlarging the constant.
For the diagonal, substitute \(k=u/\sqrt t\) and use dominated
convergence with the bound \(e^{-c|u|^2}\). The limiting integral is
\[
 \int_{\mathbb R^3}e^{-u^{\mathsf T}D_{S,J}u}
                     \frac{du}{(2\pi)^3}
 =\frac1{(4\pi)^{3/2}\sqrt{\det D_{S,J}}}.
\]
This proves \eqref{eq:ideal-site-bounds}.

Averaging a uniform site permutation makes the slot index uniform
conditional on the site. The inverse of a uniform permutation has the
same property. Every complete traversal contains one seam: if the cut
is zero, it is at the end of a forward traversal and at the start of a
reverse traversal. Free exchange transitions preserve this conditional
uniformity, since the rate to each target slot is the same and the exit
rate does not depend on the starting slot. Thus after the first seam,
the slot index is uniform conditional on the site. The site process
runs for total elapsed time \(n\beta\), proving the exact equality
in \eqref{eq:ideal-diagonal}. In particular, seam factors at a forward
and reverse junction are two fresh averages; they do not cancel as
fixed permutations would. The local limit is uniform over the finite
set \(1\le n\le2K\).

Killing restricts a path event, proving \eqref{eq:ideal-killed-upper}.
Under the interior hypothesis it cannot act before the full ideal
path leaves \(\mathfrak B_R(z)\). The ideal localization estimate
in Lemma~\ref{lem:localization}, with the same radius and at most
\(2K\) traversals, bounds this exit probability by
\(C_K e^{-c\beta}\). This proves \eqref{eq:ideal-killed-interior}.
\end{proof}

\subsection{Comparison in one exposure}

\begin{proposition}[Finite direction words]\label{prop:finite-word}
For every fixed \(K\),
\begin{equation}\label{eq:finite-word}
 \left|\mathbb E_P[(W_R(\mathbf d))_{ii}]
                 -(\widehat W_R(\mathbf d))_{ii}\right|
       =o_K(\beta^{-3/2}).
\end{equation}
The error is uniform over \(1\le n\le2K\), all direction words of
length \(n\), all cubes, deterministic cuts and slots, and all allowed
pin sets \(E\) with arbitrary finitely many extra pins. The product
inside the expectation is formed before averaging over the exposure.
\end{proposition}

\begin{proof}
\emph{Midpoint reduction.}
Split the elapsed word at time \(n\beta/2\), writing
\[
 W_R=\mathsf B\mathsf C,\qquad
 \widehat W_R=\widehat{\mathsf B}\widehat{\mathsf C}.
\]
Here \(\mathsf C\) is the starting portion. If the split is at an
instantaneous operation, assign that operation to one portion exactly
once. At a junction retain the order of the individual seam factors.
For \(\eta=\ind{\{i\}}\), put
\[
 \mu_1=\mathsf C\eta,\quad \mu_2=\mathsf B^*\eta,\qquad
 \widehat\mu_1=\widehat{\mathsf C}\eta,\quad
 \widehat\mu_2=\widehat{\mathsf B}^*\eta.
\]
After extending hole vectors by zero, all four are nonnegative masses
of total at most one on the same full slot layer. Consequently
\begin{equation}\label{eq:midpoint-pairings}
 (W_R)_{ii}=\langle\mu_1,\mu_2\rangle,\qquad
 (\widehat W_R)_{ii}=\langle\widehat\mu_1,\widehat\mu_2\rangle.
\end{equation}
The second mass is propagated backward from the terminal copy of
\(i\); it is not a separately sampled exposure. This is the role of
the matching layer in Figure~\ref{fig:midpoint}.

\begin{figure}[htbp]
\centering
\begin{tikzpicture}[>=Stealth,scale=0.95]
 \node (left) at (0,0) {$\eta$};
 \node[draw,rounded corners,fill=gray!8,minimum height=9mm,minimum width=20mm]
       (mid) at (5,0) {$\mu_1,\mu_2$};
 \node (right) at (10,0) {$\eta$};
 \draw[->,thick] (left) -- node[above] {$\mathsf C$} (mid);
 \draw[->,thick] (right) -- node[above] {$\mathsf B^*$} (mid);
 \node[below=5mm of left] {$0$};
 \node[above=3mm of mid] {common hole layer};
 \node[below=5mm of mid] {$n\beta/2$};
 \node[below=5mm of right] {$n\beta$};
\end{tikzpicture}
\caption{A diagonal return is an inner product on the midpoint layer.
Both masses use the same exposure, with fresh transitions in each
traversal. Coupling either mass to an independent full ideal path gives
an expected \(\ell^1\) error; diffusion bounds the other mass in
\(\ell^\infty\). The positions are elapsed times, not lattice sites.}
\label{fig:midpoint}
\end{figure}

Let \(G\) be the good-exposure event for \(z,R\). Each half has
elapsed length \(n\beta/2\ge\beta/2\), so
Proposition~\ref{prop:diffusion}, including its adjoint statement,
gives \(\|\mu_1\|_\infty+\|\mu_2\|_\infty\le C\beta^{-3/2}\)
on \(G\). The ideal halves obey this bound without a good-exposure
condition: each slot probability is bounded by its full site
probability in \eqref{eq:ideal-site-bounds}. A half need not contain
a seam for this upper bound. The elementary pairing inequality gives
\begin{align}
 &\left|\mathbb E_P[(W_R)_{ii}]-(\widehat W_R)_{ii}\right|\notag\\
 &\quad\le C\beta^{-3/2}\mathbb E_P
    \left[\|\mu_1-\widehat\mu_1\|_1+
          \|\mu_2-\widehat\mu_2\|_1\right]
       +\mathbb P_P(G^c).\label{eq:midpoint-reduction}
\end{align}
Indeed on \(G\) subtract the two inner products by changing their
first and second arguments successively; on \(G^c\) both diagonal
entries lie in \([0,1]\). No independence of the two halves is used.
Since \(\mathbb P_P(G^c)=o_K(\beta^{-3/2})\), it remains to prove
\begin{equation}\label{eq:coupling-l1}
 \mathbb E_P\|\mu_j-\widehat\mu_j\|_1=o_K(1),\qquad j=1,2.
\end{equation}

\medskip\noindent
\emph{Coupling conditional on the exposure.}
Both masses are generated from \(i\) along a deterministic schedule
of full or partial traversals. For \(j=2\), reverse and transpose the
suffix factors. We give the coupling for either schedule. Sample the
entire actual picture with law \(\mathbb P_P\) and expose
\(\mathcal B_P\). Independently of that entire picture, sample a
full unkilled ideal path on the infinite slot lattice, with fresh
clocks and seam permutations on each traversal. Its bounded jump rate
\(\lambda_J\) ensures that it is well defined. Conditional on the
exposure, use the same fresh clock on each slot edge whenever its
endpoints are holes in the cube. At a prescribed exposed jump, perform
the corresponding hole transport. At a seam, share the uniform site
permutation if the site has no exposed slots; otherwise use an
appropriate fresh uniform completion on its remaining slots. These
rules give exactly the hole kernel conditional on the exposure. They
also apply to reverse pieces: heat generators are symmetric, transports
are inverted, and inverse uniform bijections are uniform. Repeated
traversals share the exposure but use independent fresh transitions.

Kill both paths outside \(\Lambda^\circ\) or
\(\mathfrak B_R(z)\). Until they separate these killings agree.
There are no interaction bonds from \(\Lambda^\circ\) to the
exterior of \(\Lambda\), by the choice of boundary thickness.
A discrepancy before killing therefore requires an exposed slot at a
site whose bar-coordinate is within distance \(2B_*\) of the full
ideal path, at some traversed time. Include both sides of jumps and
instantaneous operations in this event. This covers a missing clock
edge, a forced transport affecting the common path, and a seam
completion different from a full site permutation.

\medskip\noindent
\emph{Averaged encounter bound.}
We bound the probability of this encounter unconditionally. Partition
the schedule into \(O_K(\beta)\) deterministic intervals of elapsed
length at most one, each within a single directed traversal piece;
assign endpoint operations once. Take
\(h_\beta=\lceil\log\beta\rceil\). In one such interval let
\(w\) be the bar-coordinate of the full ideal path at its initial
time. Since each jump changes bar-coordinate by at most \(B_*\),
the probability of displacement greater than \(B_*h_\beta\) during
the interval is at most
\begin{equation}\label{eq:ideal-short-tail}
 \frac{C^{h_\beta}}{h_\beta!}.
\end{equation}
For a deterministic \(w\), Corollary~\ref{cor:singleton-bound},
applied with pin set \(P\) and summed over slots, gives
\begin{equation}\label{eq:encounter-initial}
 \mathbb P_P\!\left(
  \text{an exposed slot lies over }
  \mathfrak B_{4B_*(h_\beta+1)}(w)
  \text{ at the initial slice}\right)
 \le C(h_\beta+1)^3\beta^{-3/2}.
\end{equation}
At a seam the set of sites occupied by exposed slots is the same on
both sides, so this estimate covers either endpoint convention.

If this initial event is absent but an exposed strand subsequently
visits \(\mathfrak B_{B_*(h_\beta+2)}(w)\) in the interval, trace
that strand back toward the initial slice. The distance between the
two regions forces at least \(h_\beta\) inter-site steps. Count the
first \(h_\beta\) steps from a visit. Their actual marks are distinct:
this interval lies in one traversal, and for large \(\beta\) its
length is less than \(\beta\), so unwrapped physical time is monotone
and injective modulo \(\beta\). There are at most
\(C(h_\beta+1)^3\) starting sites in the smaller region and a bounded
number of next-site and slot-edge choices per step. By the factorial
measure domination of Lemma~\ref{lem:factorial}, the probability of
such a list is at most
\begin{equation}\label{eq:encounter-entrance}
 C(h_\beta+1)^3 C^{h_\beta}
  \int_{0<s_1<\cdots<s_{h_\beta}<1}ds_1\cdots ds_{h_\beta}
 =C(h_\beta+1)^3\frac{C^{h_\beta}}{h_\beta!}.
\end{equation}
Discarding slot compatibility only enlarges the count. A seam can
change a slot index but no site, so it merely contributes a bounded
slot choice; unwrapping preserves the simplex volume. Reversing a
traversal reverses the order of the times and gives the same bound.
Thus \eqref{eq:encounter-entrance} uses factorial domination under the
full picture law, without conditioning that domination on the exposure.

The random coordinate \(w\) depends only on the independent full ideal
path. The bounds \eqref{eq:encounter-initial} and
\eqref{eq:encounter-entrance}, uniform for deterministic \(w\) and
deterministic initial times, therefore apply to this random \(w\).
A union bound over the intervals bounds the encounter probability by
\begin{equation}\label{eq:encounter-probability}
 C_K\beta(h_\beta+1)^3
       \left(\beta^{-3/2}+\frac{C^{h_\beta}}{h_\beta!}\right)
 =o_K(1).
\end{equation}
The first term is \(O_K(\beta^{-1/2}(\log\beta)^3)\); the factorial
term decays faster than every fixed inverse power of \(\beta\).
Actual marks may reappear in later traversals, but the count used
distinctness only within one interval. The union bound needs no
independence across these intervals.

Add a death state for killing. Conditional on each exposure, the
\(\ell^1\) distance between the two live output masses is at most
twice the failure probability of this coupling. Average over the
picture and use \eqref{eq:encounter-probability}. This proves
\eqref{eq:coupling-l1} for both schedules and, by
\eqref{eq:midpoint-reduction}, proves \eqref{eq:finite-word}.
\end{proof}

\subsection{Uniform return and down-spin bounds}

\begin{theorem}[One-slot bounds with the sharp coefficient]
\label{thm:return-bounds}
There is a constant \(C<\infty\), independent of \(K\), with the
following property. For every integer \(K\ge1\) there are
\(\beta_K<\infty\) and a nonnegative function
\(\epsilon_K(\beta)\to0\) as \(\beta\to\infty\), such that for
\(\beta\ge\beta_K\), every finite cube, every allowed pin set
\(E\), and every deterministic slot-time point \(i\),
\begin{equation}\label{eq:down-upper}
 d_E(i)\le\beta^{-3/2}
 \left(\frac\kappa\ell\sum_{n=1}^{2K-1}n^{-3/2}
                  +CK^{-1/2}+\epsilon_K(\beta)\right).
\end{equation}
If the site \(x\) of \(i\) additionally satisfies
\begin{equation}\label{eq:pin-free-ball}
 \mathfrak B_{R_K}(\bar x)\subseteq\Lambda^\circ,
 \qquad E\text{ has no pin over }\mathfrak B_{R_K}(\bar x),
\end{equation}
then
\begin{equation}\label{eq:down-lower}
 d_E(i)\ge\beta^{-3/2}
 \left(\frac\kappa\ell\sum_{n=1}^{K}n^{-3/2}
                                      -\epsilon_K(\beta)\right).
\end{equation}
For \(i\notin E\), the same two bounds hold for \(r(E,\{i\})\),
with the same choice of error function. The threshold and error
function are independent of the cube, cut, slot, and number, locations,
and times of all finite extra pins. Only the fixed model and geometric
choices, and \(K\), may affect them.
\end{theorem}

\begin{proof}
Suppose first that \(i\notin E\), and take the exposure law
\(\mathbb P_P\), \(P=E\cup\{i\}\). Lemma~\ref{lem:return-expansion}
and Proposition~\ref{prop:long-tail} give
\[
 r(E,\{i\})\le
 \sum_{n=1}^{2K-1}\mathbb E_P\langle\eta,Q^n\eta\rangle
 +\beta^{-3/2}\bigl(CK^{-1/2}+o_K(1)\bigr).
\]
Expand \(Q^n\) as the equally weighted average of the \(2^n\)
direction words. Replacing each word by its localized version costs
at most \(C_Ke^{-c\beta}\) in expected diagonal mass, by
Lemma~\ref{lem:localization}. Proposition~\ref{prop:finite-word} and
Lemma~\ref{lem:ideal-kernel} then imply
\[
 r(E,\{i\})\le\beta^{-3/2}
 \left(\frac\kappa\ell\sum_{n=1}^{2K-1}n^{-3/2}
                   +CK^{-1/2}+o_K(1)\right).
\]
There are only finitely many words at fixed \(K\), so all comparison
errors enter \(o_K(1)\). They do not alter the constant multiplying
\(K^{-1/2}\).

For the lower bound assume \eqref{eq:pin-free-ball}. At fixed exposure,
run the forward fresh walk with boundary and ball killing, and continue
it after visits to \(i\). Let \(A_n\) be the event that it is alive
and at \(i\) at elapsed time \(n\beta\). Before any such visit it
has stayed in the ball and hence avoided \(E\). Thus
\(q_i\ge\Pr(\bigcup_{n=1}^K A_n\mid\text{exposure})\).
The fresh transitions on successive turns give
\[
 \Pr(A_n\mid\text{exposure})=(U_{+,R}^n)_{ii},\qquad
 \Pr(A_j\cap A_n\mid\text{exposure})
     =(U_{+,R}^j)_{ii}(U_{+,R}^{n-j})_{ii}\quad(j<n).
\]
The Bonferroni lower bound therefore reads
\begin{equation}\label{eq:bonferroni}
 q_i\ge\sum_{n=1}^K(U_{+,R}^n)_{ii}
 -\sum_{1\le j<n\le K}
       (U_{+,R}^j)_{ii}(U_{+,R}^{n-j})_{ii}.
\end{equation}
On good exposures each product is \(O(\beta^{-3})\) by
Proposition~\ref{prop:diffusion}; on the complement each is at most
one, and that complement has probability \(o_K(\beta^{-3/2})\).
The expected pair sum is consequently \(o_K(\beta^{-3/2})\).
For the first sum apply Proposition~\ref{prop:finite-word} to the
all-forward word and use the interior comparison
\eqref{eq:ideal-killed-interior}. Averaging \eqref{eq:bonferroni}
gives
\[
 r(E,\{i\})\ge\beta^{-3/2}
 \left(\frac\kappa\ell\sum_{n=1}^{K}n^{-3/2}-o_K(1)\right).
\]

Finally, Proposition~\ref{prop:pin-identities} and
Corollary~\ref{cor:singleton-bound} give, uniformly,
\[
 d_E(i)=\frac{r(E,\{i\})}{1+r(E,\{i\})},\qquad
 0\le r(E,\{i\})-d_E(i)\le C\beta^{-3}.
\]
This changes only the normalized \(o_K(1)\) errors. Choose a common
threshold \(\beta_K\) and a nonnegative error function tending to
zero that dominates all these uniform error magnitudes. This proves
both estimates for \(r\) and \(d\). If \(i\in E\), then
\(d_E(i)=0\), so the upper bound still holds; the lower-bound
hypotheses already exclude this case.
\end{proof}
